% Options for packages loaded elsewhere
% Options for packages loaded elsewhere
\PassOptionsToPackage{unicode}{hyperref}
\PassOptionsToPackage{hyphens}{url}
\PassOptionsToPackage{dvipsnames,svgnames,x11names}{xcolor}
%
\documentclass[
  11pt,
  letterpaper,
  DIV=11,
  numbers=noendperiod]{scrartcl}
\usepackage{xcolor}
\usepackage[margin=1in]{geometry}
\usepackage{amsmath,amssymb}
\setcounter{secnumdepth}{-\maxdimen} % remove section numbering
\usepackage{iftex}
\ifPDFTeX
  \usepackage[T1]{fontenc}
  \usepackage[utf8]{inputenc}
  \usepackage{textcomp} % provide euro and other symbols
\else % if luatex or xetex
  \usepackage{unicode-math} % this also loads fontspec
  \defaultfontfeatures{Scale=MatchLowercase}
  \defaultfontfeatures[\rmfamily]{Ligatures=TeX,Scale=1}
\fi
\usepackage{lmodern}
\ifPDFTeX\else
  % xetex/luatex font selection
  \setmainfont[]{Calibri}
\fi
% Use upquote if available, for straight quotes in verbatim environments
\IfFileExists{upquote.sty}{\usepackage{upquote}}{}
\IfFileExists{microtype.sty}{% use microtype if available
  \usepackage[]{microtype}
  \UseMicrotypeSet[protrusion]{basicmath} % disable protrusion for tt fonts
}{}
\makeatletter
\@ifundefined{KOMAClassName}{% if non-KOMA class
  \IfFileExists{parskip.sty}{%
    \usepackage{parskip}
  }{% else
    \setlength{\parindent}{0pt}
    \setlength{\parskip}{6pt plus 2pt minus 1pt}}
}{% if KOMA class
  \KOMAoptions{parskip=half}}
\makeatother
% Make \paragraph and \subparagraph free-standing
\makeatletter
\ifx\paragraph\undefined\else
  \let\oldparagraph\paragraph
  \renewcommand{\paragraph}{
    \@ifstar
      \xxxParagraphStar
      \xxxParagraphNoStar
  }
  \newcommand{\xxxParagraphStar}[1]{\oldparagraph*{#1}\mbox{}}
  \newcommand{\xxxParagraphNoStar}[1]{\oldparagraph{#1}\mbox{}}
\fi
\ifx\subparagraph\undefined\else
  \let\oldsubparagraph\subparagraph
  \renewcommand{\subparagraph}{
    \@ifstar
      \xxxSubParagraphStar
      \xxxSubParagraphNoStar
  }
  \newcommand{\xxxSubParagraphStar}[1]{\oldsubparagraph*{#1}\mbox{}}
  \newcommand{\xxxSubParagraphNoStar}[1]{\oldsubparagraph{#1}\mbox{}}
\fi
\makeatother


\usepackage{longtable,booktabs,array}
\newcounter{none} % for unnumbered tables
\usepackage{calc} % for calculating minipage widths
% Correct order of tables after \paragraph or \subparagraph
\usepackage{etoolbox}
\makeatletter
\patchcmd\longtable{\par}{\if@noskipsec\mbox{}\fi\par}{}{}
\makeatother
% Allow footnotes in longtable head/foot
\IfFileExists{footnotehyper.sty}{\usepackage{footnotehyper}}{\usepackage{footnote}}
\makesavenoteenv{longtable}
\usepackage{graphicx}
\makeatletter
\newsavebox\pandoc@box
\newcommand*\pandocbounded[1]{% scales image to fit in text height/width
  \sbox\pandoc@box{#1}%
  \Gscale@div\@tempa{\textheight}{\dimexpr\ht\pandoc@box+\dp\pandoc@box\relax}%
  \Gscale@div\@tempb{\linewidth}{\wd\pandoc@box}%
  \ifdim\@tempb\p@<\@tempa\p@\let\@tempa\@tempb\fi% select the smaller of both
  \ifdim\@tempa\p@<\p@\scalebox{\@tempa}{\usebox\pandoc@box}%
  \else\usebox{\pandoc@box}%
  \fi%
}
% Set default figure placement to htbp
\def\fps@figure{htbp}
\makeatother





\setlength{\emergencystretch}{3em} % prevent overfull lines

\providecommand{\tightlist}{%
  \setlength{\itemsep}{0pt}\setlength{\parskip}{0pt}}



 


\KOMAoption{captions}{tableheading}
\makeatletter
\@ifpackageloaded{caption}{}{\usepackage{caption}}
\AtBeginDocument{%
\ifdefined\contentsname
  \renewcommand*\contentsname{Table of contents}
\else
  \newcommand\contentsname{Table of contents}
\fi
\ifdefined\listfigurename
  \renewcommand*\listfigurename{List of Figures}
\else
  \newcommand\listfigurename{List of Figures}
\fi
\ifdefined\listtablename
  \renewcommand*\listtablename{List of Tables}
\else
  \newcommand\listtablename{List of Tables}
\fi
\ifdefined\figurename
  \renewcommand*\figurename{Figure}
\else
  \newcommand\figurename{Figure}
\fi
\ifdefined\tablename
  \renewcommand*\tablename{Table}
\else
  \newcommand\tablename{Table}
\fi
}
\@ifpackageloaded{float}{}{\usepackage{float}}
\floatstyle{ruled}
\@ifundefined{c@chapter}{\newfloat{codelisting}{h}{lop}}{\newfloat{codelisting}{h}{lop}[chapter]}
\floatname{codelisting}{Listing}
\newcommand*\listoflistings{\listof{codelisting}{List of Listings}}
\makeatother
\makeatletter
\makeatother
\makeatletter
\@ifpackageloaded{caption}{}{\usepackage{caption}}
\@ifpackageloaded{subcaption}{}{\usepackage{subcaption}}
\makeatother
\usepackage{bookmark}
\IfFileExists{xurl.sty}{\usepackage{xurl}}{} % add URL line breaks if available
\urlstyle{same}
\hypersetup{
  pdftitle={HRM Unit Pre-Pilot Test --- Instructions},
  colorlinks=true,
  linkcolor={blue},
  filecolor={Maroon},
  citecolor={Blue},
  urlcolor={Blue},
  pdfcreator={LaTeX via pandoc}}


\title{HRM Unit Pre-Pilot Test --- Instructions}
\usepackage{etoolbox}
\makeatletter
\providecommand{\subtitle}[1]{% add subtitle to \maketitle
  \apptocmd{\@title}{\par {\large #1 \par}}{}{}
}
\makeatother
\subtitle{SCDMS --- Submission \& Commission Decision Management System}
\author{}
\date{2026-07-28}
\begin{document}
\maketitle


\emph{Version 2 --- pre-pilot dry run ahead of the wider OPSC management
pilot}

\subsection{What's different from the first
draft}\label{whats-different-from-the-first-draft}

\begin{itemize}
\tightlist
\item
  \textbf{Full coverage} --- every HR-routed submission type is in scope
  (17 total, not just one).
\item
  \textbf{Decision task allocation} --- minutes signed → task
  auto-assigned to HR Unit Manager → manager assigns to a Principal or
  Senior Officer.
\item
  \textbf{New test accounts} --- HR Unit Principal and Senior Officer,
  needed for the reassignment step, did not exist before --- created for
  this test.
\end{itemize}

\subsection{Test Accounts}\label{test-accounts}

{\def\LTcaptype{none} % do not increment counter
\begin{longtable}[]{@{}
  >{\raggedright\arraybackslash}p{(\linewidth - 4\tabcolsep) * \real{0.3333}}
  >{\raggedright\arraybackslash}p{(\linewidth - 4\tabcolsep) * \real{0.3333}}
  >{\raggedright\arraybackslash}p{(\linewidth - 4\tabcolsep) * \real{0.3333}}@{}}
\toprule\noalign{}
\begin{minipage}[b]{\linewidth}\raggedright
Role
\end{minipage} & \begin{minipage}[b]{\linewidth}\raggedright
Username
\end{minipage} & \begin{minipage}[b]{\linewidth}\raggedright
Password
\end{minipage} \\
\midrule\noalign{}
\endhead
\bottomrule\noalign{}
\endlastfoot
Ministry HR (submitter) & \texttt{hr.finance} & \texttt{Ministry123!} \\
Head of Agency / DG (endorses) & \texttt{dg.mfem} & \texttt{DG12345!} \\
PSC Officer (routes Submitted → Checklist Review) & \texttt{p.mahe} &
\texttt{Officer123!} \\
HR Unit Manager (checklist gate, allocates tasks) & \texttt{m.hrunit} &
\texttt{Manager123!} \\
HR Unit Principal (assessment, task work) --- \emph{new} &
\texttt{p.hrunit} & \texttt{Principal123!} \\
Senior Officer (alternate task assignee) --- \emph{new} &
\texttt{s.hrunit} & \texttt{Senior123!} \\
PSC Secretary (commission, minutes, decisions) & \texttt{j.iati} &
\texttt{Secretary123!} \\
Senior Admin Officer (uploads signed minutes) & \texttt{s.tari} &
\texttt{Officer123!} \\
PSC Admin (oversight, letter generation) & \texttt{admin} &
\texttt{Admin1234!} \\
\end{longtable}
}

\subsection{Phase 1 --- Smoke Test All 17 Submission Types
(breadth)}\label{phase-1-smoke-test-all-17-submission-types-breadth}

For each row below: log in as \texttt{hr.finance} → \textbf{New
Submission} → pick the agenda category shown → in the ``Specific
submission type'' dropdown, pick the exact name listed (some categories
also list older, non-digitized PSC form numbers alongside --- ignore
those, pick the named one) → fill a few realistic fields → attach any
placeholder file to each item the Required Documents panel marks
\emph{Required} → Submit to DG.

{\def\LTcaptype{none} % do not increment counter
\begin{longtable}[]{@{}
  >{\centering\arraybackslash}p{(\linewidth - 4\tabcolsep) * \real{0.3333}}
  >{\raggedright\arraybackslash}p{(\linewidth - 4\tabcolsep) * \real{0.3333}}
  >{\raggedright\arraybackslash}p{(\linewidth - 4\tabcolsep) * \real{0.3333}}@{}}
\toprule\noalign{}
\begin{minipage}[b]{\linewidth}\centering
\#
\end{minipage} & \begin{minipage}[b]{\linewidth}\raggedright
Category to pick
\end{minipage} & \begin{minipage}[b]{\linewidth}\raggedright
Specific type to select
\end{minipage} \\
\midrule\noalign{}
\endhead
\bottomrule\noalign{}
\endlastfoot
1 & Appointment / Acting Appointment & \textbf{Appointment on Probation
Submission Paper} \\
2 & Appointment / Acting Appointment & \textbf{Acting Appointment} \\
3 & Appointment / Acting Appointment & \textbf{Eligible Candidate
Notification} \\
4 & Appointment / Acting Appointment & \textbf{Unsuccessful Candidate
Notification} \\
5 & Direct Appointment / Confirmation of Appointment &
\textbf{Confirmation of Appointment Submission Paper} \\
6 & Direct Appointment / Confirmation of Appointment & \textbf{Direct
Appointment Submission Paper} \\
7 & Temporary Salaried Appointment & \textbf{Temporary Appointment
Submission Paper} \\
8 & Contract / Temporary Salaried Appointment & \textbf{Contract
Employment Submission Paper} \\
9 & Resignation / Retirement / Death & \textbf{Age Retirement
Submission} \\
10 & Resignation / Retirement / Death & \textbf{Notice of Age Retirement
Submission} \\
11 & Resignation / Retirement / Death & \textbf{Medical Retirement
Submission} \\
12 & Resignation / Retirement / Death & \textbf{Death in Service
Submission} \\
13 & Resignation / Retirement / Death & \textbf{Redundancy
Submission} \\
14 & Resignation / Retirement / Death & \textbf{Voluntary Resignation
Submission} \\
15 & Other Matters & \textbf{Secondment Submission} \\
16 & Other Matters & \textbf{Leave Payout Authorization} \\
17 & Medical Claim \emph{(auto-selected)} & \textbf{Medical Expense
Claim} \\
\end{longtable}
}

\textbf{For each one, confirm:}

\begin{itemize}
\tightlist
\item
  the digitized form fields match the real paper form;
\item
  the Required Documents list makes sense to HR;
\item
  submission is blocked until mandatory items are attached;
\item
  DG endorsement → PSC Officer routing works;
\item
  it lands in \texttt{m.hrunit}'s Manager Checklist Review queue --- and
  only there, not in VIPAM's or ODU's queue.
\end{itemize}

\emph{Log anything wrong via the feedback button as you go. Don't take
any of these past Checklist Review --- that's what Phase 2 is for.}

\subsection{Phase 2 --- Full Depth: 2 Submissions Through Decision +
Task
Allocation}\label{phase-2-full-depth-2-submissions-through-decision-task-allocation}

Pick \textbf{one Recruitment type} (suggest \#5, Confirmation of
Appointment) and \textbf{one Cessation type} (suggest \#13, Redundancy)
from Phase 1 and carry them all the way through:

\begin{enumerate}
\def\labelenumi{\arabic{enumi}.}
\tightlist
\item
  \texttt{m.hrunit} --- checklist review → assign to \texttt{p.hrunit}
\item
  \texttt{p.hrunit} --- assessment → forward to Secretary approval gate
\item
  \texttt{j.iati} --- approve gate → forward to Commission → add to a
  meeting agenda → record the Commission's decision (Approved) in the
  minutes
\item
  \texttt{s.tari} (or \texttt{admin}) --- upload the signed minutes
\item
  \textbf{Confirm automatically, without anyone touching a button:}

  \begin{itemize}
  \tightlist
  \item
    A \texttt{CommissionTask} appears in \texttt{m.hrunit}'s task
    register, already linked to the right submission.
  \item
    The submission's own ``journey'' progress bar advances to
    \textbf{Decision Entered \& Assigned} --- this is a newly-fixed
    sync, worth specifically checking it isn't stuck showing
    ``Decision'' only.
  \end{itemize}
\item
  \texttt{m.hrunit} --- from the task register, assign the task to
  \texttt{p.hrunit} for one submission and to \texttt{s.hrunit} for the
  other (testing both allowed assignee types)
\item
  \texttt{p.hrunit} / \texttt{s.hrunit} --- open their assigned task,
  update status to In Progress, then Completed
\item
  Generate the outcome letter on both submissions and check the content
  matches what was actually entered
\end{enumerate}

\subsection{What to Watch For}\label{what-to-watch-for}

\begin{itemize}
\tightlist
\item
  Does the right unit manager --- and only that manager --- see each
  submission once routed?
\item
  Does the Required Documents panel feel like a help or a hindrance to
  real HR staff --- any item that shouldn't be mandatory, or one that's
  missing?
\item
  Does the task-allocation notification actually reach
  \texttt{m.hrunit}, and the reassignment notification reach
  \texttt{p.hrunit}/\texttt{s.hrunit}?
\item
  Does the progress bar / journey view accurately reflect where things
  actually are at every stage?
\end{itemize}

Log everything through the feedback button as it happens; fill out the
System Feedback Checklist at the end (sections 2--5, 7, 8, 9 are most
relevant this round --- intake/routing, checklist review,
assignment/assessment, commission outcome, implementation tracking,
notifications).




\end{document}
